% paper2_nitrogen.tex -- Paper 2. Typeset alone with this file, or inside
% the volume through main.tex.
\documentclass[../main.tex]{subfiles}
\begin{document}
\begin{paper}{Maize Yield Response to Nitrogen Rate and Nitrogen Use Efficiency}{Nitrogen rate and nitrogen use efficiency in maize}
\label{paper:nitrogen}

\begin{paperabstract}
Nitrogen is the nutrient that most often limits cereal yields, yet a large
share of the nitrogen applied as fertiliser is not taken up by the crop. This
paper reports a nitrogen rate trial in maize with five rates of mineral
nitrogen, repeated at three sites in contrasting climatic zones. Grain yield
is related to the rate applied through a response function, from which an
agronomic optimum rate is derived, and the efficiency with which applied
nitrogen is converted to grain is expressed through standard efficiency
measures calculated against an unfertilised control. The trial data are
synthetic and serve to illustrate how rate response and nitrogen use
efficiency are calculated and reported.
\end{paperabstract}
\paperkeywords{maize; nitrogen fertiliser; rate response; agronomic efficiency; nitrogen use efficiency}

\section{Introduction}

Cereal production depends heavily on fertiliser nitrogen, and the efficiency
with which that nitrogen is used has become a central question for both
farm profitability and the environment. Reviews of cereal systems have shown
that the fraction of applied nitrogen recovered by the crop is often low, and
that the nitrogen not taken up is at risk of loss to water and air
\parencite{cassman2002}. Improving efficiency therefore depends on matching
the supply of nitrogen, in amount and timing, to the demand of the crop.

A retrospective analysis of fertiliser nitrogen in cereal production found
that recovery efficiencies measured in farmers' fields are commonly lower than
those achieved on research stations, and that the gap reflects both
management and the variable nitrogen supply of the soil itself
\parencite{ladha2005}. Nitrogen use efficiency is not a single number: it can
be expressed as the extra grain produced per unit of nitrogen applied, as the
fraction of applied nitrogen taken up by the crop, or as the grain produced
per unit of nitrogen taken up, and each measure answers a different question.

Rate trials remain the basic tool for describing how a crop responds to
nitrogen. Yield rises steeply at low rates, more slowly at intermediate rates,
and reaches a plateau once nitrogen no longer limits growth. Fitting a
response function to the yields allows the rate at which the marginal return
falls to zero to be estimated, and comparison with an unfertilised control
allows the efficiency measures to be calculated. Such trials are usually laid
out as randomised complete blocks\index{randomised complete block design}, which control for gradients in soil
fertility across the field \parencite{gomez1984}.

Fertiliser is most effective when it is used together with organic inputs and
well-adapted varieties, an approach set out as integrated soil fertility
management \parencite{vanlauwe2010}. This paper reports the mineral nitrogen
component of such a system. Its aims are to describe the yield response of
maize to five rates of nitrogen at three sites, to estimate the agronomic
optimum rate at each site, and to report the agronomic efficiency\index{nitrogen use efficiency!agronomic efficiency} of applied
nitrogen across rates and sites.

\input{\subfix{methods/paper2_methods}}

\input{\subfix{results/paper2_results}}

\input{\subfix{discussion/paper2_discussion}}

\end{paper}
\end{document}
